\subsection{A polynomial phase with logarithmic volume}\label{sec:puncturedpoly}

The puncture itself can be retained in the formal construction when
the phase is polynomial. This gives a stronger bound than the one for
$dx$: logarithmic volume does not change under a scaling of $x$.

\begin{theorem}\label{thm:puncturedpoly}
Let $K=\C((r))$, $\varphi\in K[x]$, and $h\in K\setminus\{0\}$.
On
\[
 H=K[x,x^{-1}]/(x\partial_x+x\varphi_x)K[x,x^{-1}],\qquad
 \nabla[q]=[q_r+(\varphi_r+h'/h)q],
\]
there is, after finite ramification, a full frame of rank-one formal
exponential functionals. Every zero-exponential functional in this
frame has seed value of valuation at least $v(h)$. Thus if $v(h)>0$,
the differential seed submodule has no horizontal quotient taking
$[1]$ to one. This includes a constant phase.
\end{theorem}
\begin{proof}
Write $D=\deg\varphi$, taking $D=0$ for a constant phase, including
zero. If $D\geq1$, the quotient has basis
$1,x,\ldots,x^{D-1}$. Indeed, on a negative monomial the lowest term
of $x\partial_x+x\varphi_x$ is its nonzero integer exponent times
that monomial. This eliminates all negative powers without introducing
lower ones. On polynomials, successive leading terms eliminate all
degrees at least $D$. Conversely an exact polynomial cannot come
from a Laurent polynomial with a negative lowest term, and a nonzero
polynomial preimage of degree $m$ has exact degree $m+D$.
This proves independence. For $D=0$, the quotient of $x\partial_x$
has the single basis element $1$.

Use the critical balls constructed in Theorem~\ref{thm:stationaryballs}.
At most one contains the origin. Every other ball has
$\rho_B>v(\xi_B)$, so $\gamma_B=\delta_B/\xi_B$ has positive
valuation. With the same $c_B,F_B,E_B$ and fixed functionals as before,
the local logarithmic functionals are
\begin{equation}\label{eq:puncturedaway}
 q\longmapsto h e^{c_B}\frac{\delta_B}{\xi_B}
  \ell^0_{B,j}\!\left(
      \frac{e^{E_B}q(\xi_B+\delta_B y)}{1+\gamma_B y}\right).
\end{equation}
The geometric denominator and all negative powers in $q$ expand
coefficientwise because $v(\gamma_B)>0$. Every coefficient is a
polynomial in $y$. The normalized logarithmic Jacobian is
$dx/x=\delta_B\,dy/(\xi_B+\delta_B y)$.

We retain the puncture in a separate central ball. If a critical
ball contains zero, write it as $v(x)\geq\rho_0$ and put
$\delta_0=r^{\rho_0}$. Translation within that ball gives
\[
 \varphi(\delta_0 y)=c_0+F_0(y)+E_0(r,y),\qquad
 c_0=\varphi(0),\quad E_0\in r\C[[r]][y],\quad F_0(0)=0.
\]
Here $d_0=\deg F_0\geq2$, and the ball contains $d_0-1$
critical roots with multiplicity. If no critical ball contains zero,
choose $\rho_0$ sufficiently large that the central ball is disjoint
from all critical balls and every nonconstant coefficient of
$\varphi(\delta_0 y)$ has positive valuation. Then $F_0=0$.
Such a choice exists: there are finitely many roots and coefficients.
This also handles $D\leq1$ without any critical balls.

For $d_0\geq2$, the fixed punctured quotient
\[
 H_0=\C[y,y^{-1}]/(y\partial_y+yF_0')\C[y,y^{-1}]
\]
has basis $1,y,\ldots,y^{d_0-1}$, by the same reduction argument.
For $F_0=0$, it has basis $1$ and dual functional the constant term.
For each element $\lambda_j$ of the dual basis define
\begin{equation}\label{eq:puncturedcentral}
 q\longmapsto h e^{c_0}\lambda_j
                   \!\left(e^{E_0}q(\delta_0 y)\right).
\end{equation}
There is no factor $\delta_0$: $dx/x=dy/y$. Each coefficient is
a Laurent polynomial, so the expression is well defined.

We check both identities, including parameter dependence. Spatial
exactness in the central ball is conjugation by $e^{E_0}$ of
$y\partial_y+yF_0'$. The derivative discrepancy is
$\delta_0'/\delta_0$ times the same exact expression applied to
$e^{E_0}q$. On an outside ball put $x=\xi_B+\delta_B y$ and
$w=\xi_B'+\delta_B'y$. After removing $h e^{c_B}$, the derivative
discrepancy in \eqref{eq:puncturedaway} is
\[
 e^{E_B}\left[
   \left(\frac{\delta_B'}x-\frac{\delta_B w}{x^2}\right)q
        +\frac{\delta_B w}x(q_x+\varphi_xq)\right]
 = (\partial_y+F_B')\left(e^{E_B}\frac wxq\right).
\]
It is annihilated by $\ell^0_{B,j}$. The identical product rule
proves spatial exactness there. Thus every displayed functional is
horizontal for the stated connection.

Completeness follows from polynomial separation of these disjoint
balls, not merely their number. For the central ball use the monomials
of the basis of $H_0$; for each outside ball use
$1,\ldots,y^{d_B-2}$. Lemma~\ref{lem:ballseparation} gives polynomial
amplitudes with exactly these reductions on their own ball and zero
on the others. Normalize the central columns by $h e^{c_0}$ and
the outside columns by $h e^{c_B}\delta_B/\xi_B$. Their matrix
reduces to the identity. If the central ball contains critical roots,
it contributes their multiplicity plus one; otherwise it contributes
one. Together with all outside roots this gives $D$ dimensions for
$D\geq1$, or one for $D=0$, the full quotient rank.

As in Theorem~\ref{thm:stationaryballs}, this frame is a direct sum
of rank-one exponential factors after ramification. For a zero factor,
\eqref{eq:puncturedcentral} has seed valuation at least $v(h)$;
\eqref{eq:puncturedaway} has valuation at least
$v(h)+\rho_B-v(\xi_B)>v(h)$. Semisimplicity now extends a quotient
of the differential seed submodule to the full module. Nonzero
exponential factors cannot map horizontally to the Laurent field.
A constant combination of the zero-factor seed values cannot be
one if $v(h)>0$.
\end{proof}

\begin{remark}[Why two irregular ends require another argument]
For $\varphi=r(x+x^{-1})$ and $h=1$, put $\theta=r\partial_r$.
Every logarithmic period satisfies
\[
 (\theta^2-4r^2)I=0.
\]
Indeed $y=x+x^{-1}$ and $z=x-x^{-1}$ obey
$(x\partial_x+x\varphi_x)z=y+rz^2$ and $y^2-z^2=4$.
These identities give the equation by integration by parts.
The quotient has rank two, and $[1],[y]$ are independent:
its basis is $x^{-1},1$ and $[y]=2[x^{-1}]$.

An explicit formal basis of scalar solutions is
\[
 I=\sum_{n\geq0}\frac{r^{2n}}{(n!)^2},\qquad
 J=I\log r-\sum_{n\geq1}\frac{H_n r^{2n}}{(n!)^2},
 \quad H_n=\sum_{j=1}^n\frac1j.
\]
The coefficient recurrence $n^2a_n=a_{n-1}$ proves the first
equation; $H_n-H_{n-1}=1/n$ proves the second.
Their $\theta$-Wronskian is one: its derivative is zero by the
equation and its constant term is one. Thus the logarithm is a full
rank-two phenomenon, with nontrivial unipotent monodromy that no
finite ramification removes. These are the familiar modified Bessel
solutions; see~\cite[10.25.1, 10.31.2]{DLMF}.
This example does not refute the desired logarithmic bound:
multiplication by $h=r$ makes both seed values tend to zero.
It does refute an extension of the semisimple polynomial-frame theorem
to all Laurent phases. The remaining proof must allow such extensions.
\end{remark}

The residue dimension at the puncture is essential. Dropping it would
give rank $D-1$ instead of $D$ and an incomplete frame. This theorem
still permits nonconvergent formal coefficients, but is a valued-field
statement. It does not cover a phase with both positive and negative
powers of $x$.

\begin{corollary}\label{cor:oneend}
Suppose $\varphi(t,x)$ is rational in $t$ and polynomial either in
$x$ or in $x^{-1}$. For the logarithmic-volume connection normalized
by $h=1/t$, no horizontal formal functional at another point can have
seed value one.
\end{corollary}
\begin{proof}
For a polynomial in $x$, apply Theorem~\ref{thm:puncturedpoly} at
$r=1/t=0$, where $v(h)=1$. The rational seed restriction is justified
as in Corollary~\ref{cor:leadingobstruction}. For a polynomial in
$x^{-1}$ use $y=1/x$; $dx/x=-dy/y$ only changes a constant factor.
\end{proof}
